\subsection{Управление скоростью: cmd\_vel}
\label{sec:ros2-cmdvel}

Цель раздела: управлять дроном скоростью, публикуя сообщения в топик, и держать для этого ровный ритм таймером.

Скорость задаётся не сервисом, а \emph{публикацией} в топик \fpath{/geoscan/flight/cmd_vel}. Тип сообщения --- \texttt{geometry\_msgs/\allowbreak TwistStamped}: заголовок и \emph{twist} (линейная и угловая скорость).

\begin{center}
\small
\begin{tabularx}{\linewidth}{@{}>{\raggedright\arraybackslash}p{0.34\linewidth} X@{}}
\toprule
\textbf{Поле} & \textbf{Смысл} \\
\midrule
\verb|header.frame_id| & система осей: \verb|'base_link'| --- оси коптера, \verb|'map'| --- мировые; пустое или другое значение сообщение отключает (предупреждение в журнале) \\
\verb|twist.linear.x, y, z| & линейная скорость, м/с \\
\verb|twist.angular.z| & скорость поворота, рад/с; предел $\pm 1.82$ \\
\verb|twist.angular.x, y| & игнорируются \\
\bottomrule
\end{tabularx}
\end{center}

Топик работает, только если в конфиге включён \verb|topics.cmd_vel| (в поставляемых конфигах включён). Пример квадрата скоростью:

\begin{codefile}[python]{python-onboard/ros2/ros2_03_cmd_vel_square.py}
\end{codefile}

\paragraph{Что происходит.} Список \verb|LEGS| --- четыре стороны квадрата как пары скоростей $(v_x, v_y)$ в осях \verb|map|: на каждую уходит \verb|LEG_TICKS| = 25 тиков таймера по 0.1~с, то есть $25 \times 0.1 \times 0.4 = 1$~м. Таймер публикует скорость текущей стороны и заканчивает работу флагом \verb|finished|. Сначала таймер отменён (\verb|cancel()|): пока дрон не в воздухе, публиковать нечего; после взлёта его запускают \verb|reset()|. Главный цикл только ждёт флаг. Конец --- команда \verb|land| и ожидание \verb|fly_state = 0|.

\begin{rulebox}[Скорость нужно подтверждать постоянно]
Команда \verb|cmd_vel| не «включает» скорость навсегда. В симуляторе каждое сообщение держит скорость 0.5~с, затем дрон зависает; на земле команда игнорируется с предупреждением. Поэтому публикуйте скорость таймером раз в 0.1~с --- и дрон сам остановится, если программа зависнет. (У \texttt{pioneer\_sdk2} то же делает аргумент \verb|interval|.) На настоящем дроне срок удержания сверьте опытом.
\end{rulebox}

Перед \verb|land| после \verb|cmd_vel| выдержите короткую паузу (\verb|spin_once| несколько раз): запоздавшее сообщение скорости не должно прийти поверх посадки. Квадрат из скоростей замыкается хуже квадрата из точек: ошибка накапливается так же, как у относительных перемещений в главе~\ref{sec:sdk2}.

\begin{codebox}[Скорость из консоли]
\small
\begin{verbatim}
ros2 topic pub -r 10 /geoscan/flight/cmd_vel \
  geometry_msgs/msg/TwistStamped \
  "{header: {frame_id: 'base_link'},
    twist: {linear: {x: 0.5, y: 0.0, z: 0.0}, angular: {z: 0.0}}}"
\end{verbatim}
\end{codebox}

\begin{infobox}[Ограничение курса и линейной скорости]
Предел $\pm 1.82$ рад/с относится к \verb|angular.z|: в симуляторе он понят как \emph{ограничение} (значение обрезается), тогда как в \texttt{pioneer\_sdk2} команду с большей скоростью поворота отбрасывают целиком. Линейные скорости эмуляция не ограничивает: предел задаёт автопилот.
\end{infobox}

\begin{taskbox}[Квадрат скоростью в base\_link]
Переделайте пример: \verb|frame_id = 'base_link'|, сторона --- полёт вперёд ($v_x = 0.4$), затем поворот на месте $\omega_z$ около $0.8$~рад/с в течение $\pi/2 / 0.8 \approx 2$~с. Сделайте замкнутый квадрат и сравните его с квадратом в \verb|map|.
\expected{дрон проходит четыре стороны и поворачивается налево на каждой вершине (положительная скорость поворота --- влево); возвращение к старту хуже, чем у квадрата в \texttt{map}, за счёт накопления ошибки поворота.}
\mission{Мини~2 SDK2}{47, 48}
\end{taskbox}

\subsection{Виртуальный пульт: rc/set\_channels}

Топик \fpath{/geoscan/rc/set_channels} (\texttt{gs\_interfaces/\allowbreak RcChannels}) подаёт на автопилот значения каналов пульта --- как будто они пришли с настоящего пульта: поля \verb|channel_1| \ldots \verb|channel_8|. Тот же приём есть и в \texttt{pioneer\_sdk2} (виртуальный пульт). Справочник Геоскана требует публиковать каналы \emph{постоянно}, чтобы удерживать связь. Топик создаётся при \verb|topics.rc: true|.

\begin{codefile}[python]{python-onboard/ros2/ros2_09_rc_channels.py}
\end{codefile}

\paragraph{Что происходит.} Таймер \verb|create_timer(0.1, self.send)| раз в 0.1~с публикует все восемь каналов: ручки по центру (1500), восьмой канал --- тумблер (2000 --- включён). В симуляторе значения --- как у \verb|pioneer_send_rc_channels|: от 1000 до 2000, центр 1500. Подписка на \fpath{rc/trigger} печатает значение канала, когда оно меняется (в эмуляции --- 8-й канал). Через несколько секунд программа завершается: публикация прекращается.

\begin{warnbox}[RC в воздухе]
Каналы, посланные из программы, приходят на тот же путь ручного управления, что и пульт. Не подавайте их, пока не знаете, что делает каждый канал вашего набора параметров (\texttt{Copter\_man\_rcMode0..2}, раздел~\ref{sec:ros2-params}), и не рассчитывайте на них как на замену аварийной остановки с пульта.
\end{warnbox}

\begin{taskbox}[Тумблер и отклик]
Опубликуйте восьмой канал в положении 1000 три секунды, затем 2000 ещё три и после этого прекратите публикацию. Печатайте значения \fpath{rc/trigger} и убедитесь, что они меняются вслед за каналом.
\expected{в журнале после смены канала приходит новое значение (1000, затем 2000); когда публикация прекращена, новых сообщений нет.}
\mission{Мини~2 SDK2}{49}
\end{taskbox}
